\section{r81 Matter-frame epsilon terminality and Planck/information scale split}
\label{sec:r81-matter-frame-epsilon-terminality}

This source-only addendum records the physical interpretation suggested by the r60--r80
endpoint.  The terminal blocker
\[
0_{\rm termdata}\not\Rightarrow 0_{\rm all}
\]
is not merely a missing algebraic shortcut.  It is the correct typing if the matter
readout is an operational, finite-resolution readout rather than an exact all-readout
zero object.

\subsection*{Scale typing}
\label{subsec:r81-scale-typing}

Let \(\mathsf I\) denote the information-side layer, \(\mathsf M\) the matter/detector
readout, and \(\mathsf P\) the photon/clock-boundary readout.  The matter readout is
modelled as a coarse or terminal projection
\[
\pi_{\mathsf I\to\mathsf M}:\mathsf I\longrightarrow \mathsf M.
\]
It is therefore not required that a native information-scale refinement appears as a
shorter material length.  The frame-typed scale relation is
\[
\ell_{\mathsf I}^{\rm native}\leq \ell_{\mathsf M}^{\rm eff},
\]
where \(\ell_{\mathsf M}^{\rm eff}\) is the matter-frame terminal resolution.  In a
Planck-compatible normalization one may set
\[
\varepsilon_{\mathsf M}^{\rm len}=\frac{\ell_{\rm P}}{L_{\rm ref}}
\]
or, in the already dimensionless detector chain,
\[
\theta_{\rm det}=\frac1{256}.
\]
These are terminal scales, not exact all-readout zeros.

\subsection*{Matter-frame epsilon terminality theoremlet}
\label{subsec:r81-matter-epsilon-theoremlet}

\paragraph{Theoremlet r81.}
Assume a matter-frame residual \(\Phi\) is read only through a terminal matter detector
\(D_{\varepsilon_{\mathsf M}}\).  If
\[
\Phi\leq \varepsilon_{\mathsf M},
\]
then
\[
D_{\varepsilon_{\mathsf M}}(\Phi)=0_{\mathsf M}^{\varepsilon}
\quad\Rightarrow\quad
0_{\rm termdata}.
\]
However, in the guarded grammar,
\[
0_{\mathsf M}^{\varepsilon}\not\Rightarrow 0_{\rm all},
\qquad
0_{\rm termdata}\not\Rightarrow 0_{\rm all}.
\]
Consequently matter-frame closure is an effective/epsilon closure, not exact global
all-readout closure.

\paragraph{Proof sketch.}
The detector map is typed as a terminal projection.  Its zero is a record zero: it states
that the residual has fallen below the matter-frame operational terminal scale.  The r63
and r64 ledgers already show this for \(\theta_{\rm det}=1/256\), where
\[
\Phi_5^{(20)}\leq \frac1{256}
\]
while \(\Phi_5^{(n)}>0\) for every finite \(n\).  The r73--r80 bridge and quotient scans
show that no guard-admissible implication from terminal data to exact all-readout zero is
available.  Therefore the implication stops at terminal data. \(\square\)

\subsection*{Photon-frame null note}
\label{subsec:r81-photon-frame-null-note}

The theoremlet does not prove that the photon/clock frame carries an exact global zero.
It only permits a cleaner typing:
\[
0_{\mathsf M}^{\varepsilon}\neq 0_{\rm all},
\qquad
0_{\mathsf P}\text{ remains a separate candidate null object.}
\]
Thus a photon-frame or boundary-clock zero, if later constructed, must still be related to
\(0_{\rm all}\) by a separate theorem.  It must not be imported from the matter-frame
terminal detector.

\subsection*{Effective closure layer}
\label{subsec:r81-effective-closure-layer}

We introduce the scoped closure layer
\[
\mathsf{Closure}_{\mathsf M,\rm eff}^{\varepsilon}
\]
with the admissibility rule
\[
\Phi\leq\varepsilon_{\mathsf M}
\quad\Longrightarrow\quad
0_{\mathsf M}^{\varepsilon}\quad\Longrightarrow\quad 0_{\rm termdata},
\]
and the explicit non-promotion guard
\[
\mathsf{Closure}_{\mathsf M,\rm eff}^{\varepsilon}
\not\Rightarrow
\mathsf{Closure}_{\rm all}^{\rm glob}.
\]
This separates three levels:
\[
\mathsf{Closure}_{\rm loc},
\qquad
\mathsf{Closure}_{\mathsf M,\rm eff}^{\varepsilon},
\qquad
\mathsf{Closure}_{\rm all}^{\rm glob}.
\]
The first two are compatible with the v4.7.0 guarded release scope.  The third remains a
forbidden claim.

\subsection*{r81 guard update}
\label{subsec:r81-guard-update}

The r81 scale-typing lift changes the interpretation of the terminal blocker but does not
spend any global guard:
\[
R_{\mathsf M,\rm eff}^{r81}=(0,0,0),
\qquad
R_{\rm glob}^{(r81,+)}=(0,1,1,1,1),
\]
with
\[
\Gamma_{\rm guard}^{r81}=(1,1,1,1,1).
\]
The release wording may now say controlled matter-frame/effective closure.  It must not
say exact finite \(C5\) zero closure, exact all-readout terminality, \(C6/T2\), global
photon-clock/RH theorem, or global all-readout closure.


% -----------------------------------------------------------------------------
\section*{v4.7.0-r82: Tachyon-frame signed-scale separation}
\label{sec:r82-tachyon-frame-signed-scale}

\paragraph{Purpose.}
The r81 matter-frame theoremlet separates effective matter-frame terminality from exact
all-readout zero.  The present r82 note adds the corresponding tachyon/information-facing
scale typing.  It answers a possible ambiguity in the phrase ``negative length'': the
framework does not introduce negative material lengths.  It introduces a signed or
spacelike norm class for a tachyon-frame readout, relative to the matter-frame metric
convention.

\subsection*{Signed scale convention}
\label{subsec:r82-signed-scale-convention}

Let \(\mathsf T\) denote a tachyonic, superluminal, or information-facing readout channel.
Its native scale is written as a positive magnitude
\[
\lambda_{\mathsf T}>0.
\]
When this channel is expressed through the matter-frame interval convention, it is typed as
opposite-sign/spacelike:
\[
\ell_{\mathsf T}^{2}=-\lambda_{\mathsf T}^{2}
\]
for the matter-frame signature used in this branch.  If the opposite metric signature is
used, the sign convention is reversed, but the content is unchanged: the tachyon-frame
scale carries the opposite interval sign from an ordinary matter-frame material length.

The guard is
\[
\ell_{\mathsf T}^{2}<0
\not\Rightarrow
L_{\mathsf M}<0.
\]
Thus the negative sign is not a negative material length.  It is a signature marker for a
spacelike, dual, or information-facing scale class.

\subsection*{Theoremlet r82: tachyon-frame signed-scale separation}
\label{subsec:r82-theoremlet-signed-scale-separation}

\paragraph{Theoremlet.}
Assume the r81 matter-frame effective terminal layer
\(\mathsf{Closure}_{\mathsf M,\rm eff}^{\varepsilon}\) and the r60+ guarded relative-seam
extension.  A tachyon-frame native scale may be finer than the matter-frame effective
terminal scale,
\[
\lambda_{\mathsf T}\leq \ell_{\mathsf M}^{\rm eff},
\]
without becoming a shorter material length in the matter frame.  Matter-frame access is
mediated by a projection or coarse-graining map
\[
\pi_{\mathsf T\to\mathsf M}:\mathsf T\to\mathsf M,
\]
so that tachyon-frame substructure is read in the matter frame only as projected residual,
phase/address data, or terminal data.  Consequently
\[
0_{\mathsf T}^{\varepsilon}
\not\Rightarrow
0_{\mathsf M}^{\varepsilon}
\not\Rightarrow
0_{\rm all},
\]
and
\[
\ell_{\mathsf T}^{2}<0
\not\Rightarrow
L_{\mathsf M}<0.
\]

\paragraph{Interpretation.}
The tachyon frame is allowed to carry a native information or phase resolution below the
matter-frame terminal scale.  The matter frame does not measure this as a material length
shorter than its effective Planck-compatible terminal scale.  It sees it through
coarse-graining, residual phase structure, detector records, or relay data.

\paragraph{Proof sketch.}
The r81 theoremlet already establishes that matter-frame residuals below
\(\varepsilon_{\mathsf M}\) yield \(0_{\mathsf M}^{\varepsilon}\) and terminal data, but
not exact all-readout zero.  A tachyon-frame scale has different variance and interval
signature from a matter-frame material length.  The comparison therefore requires an
explicit typed comparison map or projection.  The projection
\(\pi_{\mathsf T\to\mathsf M}\) can preserve relay-relevant terminal data, but it does not
convert an opposite-sign interval into a negative matter-frame length and does not provide
an all-readout zero bridge.  The implication therefore stops at projected terminal or
residual data. \(\square\)

\subsection*{Frame-scale table}
\label{subsec:r82-frame-scale-table}

The branch now distinguishes four scale/readout roles:
\[
\begin{array}{c|c|c}
\text{readout} & \text{scale type} & \text{zero status}\\
\hline
\mathsf M & \ell_{\mathsf M}^{\rm eff}\sim\varepsilon_{\mathsf M} &
0_{\mathsf M}^{\varepsilon}\Rightarrow0_{\rm termdata},\;0_{\mathsf M}^{\varepsilon}\not\Rightarrow0_{\rm all}\\
\gamma & \ell_{\gamma}\text{ / clock selector} &
0_{\gamma}\text{ remains a separate candidate null object}\\
\mathsf T & \ell_{\mathsf T}^{2}=-\lambda_{\mathsf T}^{2},\;\lambda_{\mathsf T}>0 &
0_{\mathsf T}^{\varepsilon}\not\Rightarrow0_{\rm all}\\
\mathsf I & \ell_{\mathsf I}^{\rm native}\text{ / address scale} &
\text{native information scale, not a material length}
\end{array}
\]

This prevents three overclaims:
\[
\lambda_{\mathsf T}<\ell_{\mathsf M}^{\rm eff}
\not\Rightarrow
\text{shorter matter length},
\]
\[
\ell_{\mathsf T}^{2}<0
\not\Rightarrow
\text{negative material length},
\]
\[
0_{\mathsf T}^{\varepsilon}
\not\Rightarrow
0_{\rm all}.
\]

\subsection*{r82 guard update}
\label{subsec:r82-guard-update}

The r82 signed-scale typing adds a tachyon-frame separation layer but does not spend any
global guard:
\[
R_{\mathsf T,\rm scale}^{r82}=(0,0,0),
\qquad
R_{\rm glob}^{(r82,+)}=(0,1,1,1,1),
\]
with
\[
\Gamma_{\rm guard}^{r82}=(1,1,1,1,1).
\]
The effective matter-frame residual remains zero only inside the scoped
\(\mathsf{Closure}_{\mathsf M,\rm eff}^{\varepsilon}\) layer.  The tachyon-frame signed
scale does not imply exact finite \(C5\) zero closure, exact all-readout terminality,
\(C6/T2\), a global photon-clock/RH theorem, or global all-readout closure.



% ============================================================================
\section*{v4.7.0-r83 Source-Only Addendum: Effective/Local Closure Reclassification}
\label{sec:r83-effective-local-closure-reclassification}
% ============================================================================

\paragraph{Purpose.}
The r81--r82 frame-scale separation turns several previously open-looking residuals into
scoped closure statements.  This r83 addendum does not claim global all-readout closure.
Instead it reclassifies the status of the affected objects by separating effective/local
closure from exact/global closure.

\subsection*{r83 closure classes}
\label{subsec:r83-closure-classes}

We use the following scoped closure labels:
\[
\mathsf{Closed}_{\rm eff/local},
\qquad
\mathsf{Closed}_{\rm typed},
\qquad
\mathsf{Open}_{\rm exact/global}.
\]
Here \(\mathsf{Closed}_{\rm eff/local}\) means closed only inside the declared
matter-frame effective, detector-terminal, or local relay/clock layer.  It is not an
all-readout zero claim.  \(\mathsf{Closed}_{\rm typed}\) means that an ambiguity or
misreading has been eliminated by a typing/separation guard.  \(\mathsf{Open}_{\rm
exact/global}\) means that no exact all-readout, finite-zero, or T2/global witness has
been supplied.

\subsection*{r83 reclassification theoremlet}
\label{subsec:r83-reclassification-theoremlet}

\paragraph{Theoremlet r83.}
Under the r60 conservative relative-seam extension and the r81--r82 frame-scale guards,
the following reclassification is admissible:
\[
\begin{array}{c|c|c}
\text{object} & \text{r83 status} & \text{scope guard}\\
\hline
\text{matter-frame zero problem} & \mathsf{Closed}_{\rm eff/local} &
0_{\mathsf M}^{\varepsilon}\Rightarrow0_{\rm termdata},\;0_{\mathsf M}^{\varepsilon}\not\Rightarrow0_{\rm all}\\
\text{tachyon negative-length ambiguity} & \mathsf{Closed}_{\rm typed} &
\ell_{\mathsf T}^{2}<0\not\Rightarrow L_{\mathsf M}<0\\
OP5_{\rm eff} & \mathsf{Closed}_{\rm eff/local} &
\Phi_5^{(20)}\le 1/256\Rightarrow0_{\mathsf M}^{\varepsilon}\\
OP5_{\rm exact} & \mathsf{Open}_{\rm exact/global} &
\Phi_5^{(n)}>0\text{ for finite }n\\
C5_{\rm det/eff} & \mathsf{Closed}_{\rm eff/local} &
D_{1/256}(\Phi_5)=0_{\rm det}\text{ at }n=20\\
C5_{\rm zero} & \mathsf{Open}_{\rm exact/global} &
\nexists N<\infty:\Phi_5^{(N)}=0\\
W_{\rm clk}^{\rm loc} & \mathsf{Closed}_{\rm eff/local}\text{ conditionally} &
\Delta_{\gamma}<1/5\\
W_{\gamma}^{\rm glob} & \mathsf{Open}_{\rm exact/global} &
\text{no global photon-clock/RH theorem}\\
\mathsf C_{\rm mix}^{\rm typed} & \mathsf{Closed}_{\rm typed} &
\text{local typed mixed comparison only}\\
C6/T2 & \mathsf{Open}_{\rm exact/global} &
0_{\rm termdata}\not\Rightarrow0_{\rm all}
\end{array}
\]

\paragraph{Proof sketch.}
The matter-frame implication is exactly the r81 effective terminality guard:
\[
\Phi\le\varepsilon_{\mathsf M}\Rightarrow0_{\mathsf M}^{\varepsilon}\Rightarrow0_{\rm termdata},
\qquad
0_{\mathsf M}^{\varepsilon}\not\Rightarrow0_{\rm all}.
\]
The tachyon-frame ambiguity is removed by the r82 signed-scale convention
\(\ell_{\mathsf T}^{2}=-\lambda_{\mathsf T}^{2}\), \(\lambda_{\mathsf T}>0\), with the
explicit guard \(\ell_{\mathsf T}^{2}<0\not\Rightarrow L_{\mathsf M}<0\).  The OP5/C5
finite detector statement follows from the previously established threshold
\(\Phi_5^{(20)}\le1/256\).  The exact OP5/C5 statement remains open because the bounded
collar recurrence remains positive for every finite step.  The clock and mixed-comparison
closures inherit their local/typed guards from r65 and r69 respectively.  The final T2/global
row remains open because r75--r80 freeze the terminal blocker
\(0_{\rm termdata}\not\Rightarrow0_{\rm all}\).  \(\square\)

\subsection*{r83 status ledger}
\label{subsec:r83-status-ledger}

The reclassification changes the narrative, not the global residual.  The correct public
reading is:
\[
\text{effective/local/frame-typed closures gained; exact/global closure not claimed.}
\]
Thus the effective/local ledger is
\[
R_{\rm eff/local}^{r83}=(0,0,0,0,0),
\]
for the scoped classes
\[
(0_{\mathsf M}^{\varepsilon},\;\ell_{\mathsf T}^{2},\;OP5_{\rm eff},\;C5_{\rm det/eff},\;W_{\rm clk}^{\rm loc}),
\]
while the exact/global ledger remains
\[
R_{\rm exact/global}^{r83}=(1,1,1,1,1),
\]
for
\[
(C5_{\rm zero},\;0_{\rm all},\;C6/T2,\;W_{\gamma}^{\rm glob},\;W_{\rm all}).
\]
Equivalently, the main global residual remains
\[
R_{\rm glob}^{(r83,+)}=(0,1,1,1,1),
\qquad
\Gamma_{\rm guard}^{r83}=(1,1,1,1,1).
\]
The first component stays zero only in the r60+ relative-seam setting; the remaining
components are not promoted by the r83 reclassification.

\subsection*{r83 release-scope implication}
\label{subsec:r83-release-scope-implication}

For v4.7.0 release language, r83 supports the phrase ``controlled effective/local
closure'' but forbids replacing it by ``global closure.''  The safe formulation is:
\[
\boxed{\text{v4.7.0 establishes guarded effective/local and frame-typed closures, not T2/global closure.}}
\]
The forbidden formulations remain:
\[
C5=0,
\qquad
0_{\rm termdata}\Rightarrow0_{\rm all},
\qquad
C6/T2,
\qquad
W_{\gamma}^{\rm glob},
\qquad
W_{\rm all}.
\]


\clearpage
\section*{v4.7.0-r84 Post-r83 Full-Monolith Integration and Release-Scope Refresh}
\addcontentsline{toc}{section}{v4.7.0-r84 Post-r83 Full-Monolith Integration and Release-Scope Refresh}
\label{sec:r84-post-r83-full-monolith-release-scope-refresh}

\noindent\textbf{Source-only status.}
This r84 layer integrates the r83 effective/local closure reclassification back into the
full-monolith source path after the pre-r81--r83 handoff.  It does not create a public release and it
does not create a new PDF.  The intended reserved v4.7.0 DOI remains
\href{https://doi.org/10.5281/zenodo.20232911}{10.5281/zenodo.20232911}; the previous
baseline DOI \href{https://doi.org/10.5281/zenodo.20117361}{10.5281/zenodo.20117361} remains
provenance only.

\subsection*{r84 integration ledger}
\label{subsec:r84-integration-ledger}
The full-monolith integration keeps the r81--r83 additions in the proof spine while preserving
the corrected compact title-page discipline introduced during full-monolith recovery.  In particular,
the title page carries one version DOI line only, and release notices, previous DOI provenance,
and claim-scope explanations remain in the body/provenance layers rather than on the title page.

The post-r83 source object is
\[
\mathfrak M_{r84}
=\bigl(\mathfrak M_{r83},\mathsf{Scope}_{\rm eff/local}^{r83},
\mathsf{Guard}_{\rm exact/global}^{r83},\mathsf{Meta}_{\rm DOI}^{r84}\bigr).
\]
The integration residual is
\[
R_{\rm int}^{r84}=(0,0,0,0),
\]
corresponding to source continuity, title-page compactness, DOI separation, and claim-scope
preservation.

\subsection*{r84 release-scope refresh}
\label{subsec:r84-release-scope-refresh}
The safe v4.7.0 release phrase after r83 is:
\[
\boxed{\text{guarded effective/local and frame-typed closures, not exact global closure.}}
\]
The allowed local/effective assertions are
\[
\begin{gathered}
0_{\mathsf M}^{\varepsilon}\Rightarrow0_{\rm termdata},\qquad
\ell_{\mathsf T}^{2}=-\lambda_{\mathsf T}^{2},\;\lambda_{\mathsf T}>0,\\
OP5_{\rm eff},\qquad C5_{\rm det/eff}\text{ at }\theta=1/256\text{ after }n=20,\\
W_{\rm clk}^{\rm loc}\text{ under }\Delta_{\gamma}<1/5,
\qquad \mathsf C_{\rm mix}^{\rm typed}\text{ locally.}
\end{gathered}
\]
The forbidden exact/global assertions remain
\[
C5=0,
\qquad
0_{\rm termdata}\Rightarrow0_{\rm all},
\qquad
C6/T2,
\qquad
W_{\gamma}^{\rm glob},
\qquad
W_{\rm all}.
\]
Thus the refreshed ledgers are
\[
R_{\rm eff/local}^{r84}=(0,0,0,0,0),
\qquad
R_{\rm exact/global}^{r84}=(1,1,1,1,1),
\]
and the global residual remains
\[
R_{\rm glob}^{(r84,+)}=(0,1,1,1,1),
\qquad
\Gamma_{\rm guard}^{r84}=(1,1,1,1,1).
\]

\subsection*{r84 release-path implication}
\label{subsec:r84-release-path-implication}
Because r83 changes the conceptual framing from ``unresolved failure'' to scoped
effective/local closure, any later public-release package should be regenerated from r84 or later,
not from a pre-r81--r83 handoff.  A later release candidate may be built by applying the compact-titlepage
rule and final-handoff naming discipline to this r84 source.  No such public release is
authorized by r84 itself.


